%
% frame-visibility-lifecycle.tex
%
% A Z specification of the lux Display's frame-visibility lifecycle: which
% events may move a frame between on-screen, docked and closed, and which
% events may move a scene id between unseen and known.  Formalises the R8
% rule of DES-065 ("show() is a notification, not a window-raise") together
% with the refinement bead lux-mxvy.8 adds: closed is a visibility state the
% Display owns, exactly like minimized, and a content push may not touch it.
%
% Source of record:
%   src/punt_lux/display/replica/scene_replica.py -- upsert_scene_in_frame,
%                                    close_frame, dismiss_framed_scene
%   src/punt_lux/display/replica/frame_book.py    -- ensure, minimize,
%                                    pop_frame, forget_scene, request_focus
%   src/punt_lux/display/replica/frame.py         -- Frame.minimized
%   src/punt_lux/display/frame_commands.py        -- FrameCommands.raise_it
%
% Type-check:  fuzz -t docs/frame-visibility-lifecycle.tex
% Model-check (see Section "Verification" for the exact goal predicates; every
% goal below must report NOT found):
%   probcli docs/frame-visibility-lifecycle.tex -model_check \
%       -p DEFAULT_SETSIZE 2
%   probcli docs/frame-visibility-lifecycle.tex -model_check -nodead \
%       -goal "autoFocused /= {}" -p DEFAULT_SETSIZE 2
%   probcli docs/frame-visibility-lifecycle.tex -model_check -nodead \
%       -goal "tabStolen /= {}" -p DEFAULT_SETSIZE 2
%   probcli docs/frame-visibility-lifecycle.tex -model_check -nodead \
%       -goal "userClosed /\ frames /= userClosed" -p DEFAULT_SETSIZE 2
%   probcli docs/frame-visibility-lifecycle.tex -model_check -nodead \
%       -goal "card({f|f:FRAME & f:frames & f:userClosed & vis(f)=vOpen})>0" \
%       -p DEFAULT_SETSIZE 2
%
% The fidelity control that must report each of those four goals FOUND is
% docs/frame-visibility-lifecycle_buggy.tex.
%
% Formatting follows the fuzz house rules: \quad~ for continuation indent
% (never \t1..\t9), schema lines under 80 columns, predicates broken at
% \land/\lor/\implies with the operator starting the continuation line.
%

\documentclass[a4paper,10pt,fleqn]{article}
\usepackage[margin=1in]{geometry}
\usepackage{fuzz}
\usepackage[colorlinks=true,linkcolor=blue,citecolor=blue,urlcolor=blue]{hyperref}

\begin{document}

\title{The lux Frame-Visibility Lifecycle: a Z Specification}
\author{Formal model of the content/visibility separation behind DES-065 R8}
\date{August 2026}
\maketitle

% ============================================================
\section{Introduction}
% ============================================================

Two different authorities write to the same frame.  A \emph{client} writes
its content --- \texttt{show()} installs a scene, \texttt{update()} replaces
one, an empty push or a manifest purge takes one away.  The \emph{user}
writes its visibility --- the collapse button docks a frame, the close
button dismisses it, a dock pill or an applet's menu entry brings it back.
DES-065's R8 states the rule these two must obey: an agent's push is a
notification, not a window-raise.  Bead \texttt{lux-mxvy.8} adds the
refinement this model exists to pin down --- \emph{closed} is a visibility
state the Display owns, in the same category as \emph{minimized}, and not
the erasure of the frame it is today.

Two shipped defects, both found downstream in vox, are the same rule broken
in opposite directions.

\begin{description}
  \item[Bug A (vox-640w).]  A frame the user put away cannot be brought
    back.  Today \texttt{close\_frame} removes the frame outright, so the
    one gesture that restores a frame --- \texttt{raise\_frame} --- has
    nothing left to act on and answers \texttt{raised: false}.  A frame the
    user closed is only ever restored by accident, when the next content
    push happens to recreate it.
  \item[Bug B (vox-7l2d).]  A frame the user closed reopens on its own the
    next time its content changes in the background.  Today
    \texttt{close\_frame} calls \texttt{forget\_scene}, so the scene id
    returns to \emph{unseen}; the next push therefore reads as new, and
    \texttt{upsert\_scene\_in\_frame}'s \texttt{is\_new} branch clears
    \texttt{minimized} and requests focus.  A track change reopens a window
    the user shut.
\end{description}

They share one cause.  \texttt{is\_new} --- ``this scene id is not in this
frame'' --- is a fact about \emph{content}, and it is wired to affordances
that belong to \emph{visibility}.  Because a close erases content, the user's
visibility decision is fed back into the content axis and destroyed there.

This document models the two axes as independent state and states the
separation as invariants, so that ProB can enumerate the reachable states and
either confirm the separation holds or produce the exact offending trace.  The
companion \texttt{frame-visibility-lifecycle\_buggy.tex} is the same model
with today's behaviour restored in exactly two operations, and it must reach
the states this one cannot.

% ============================================================
\section{Basic Types}
% ============================================================

\begin{zed}
  [FRAME, SCENE]
\end{zed}

\texttt{FRAME} is the set of frame identifiers a client may name;
\texttt{SCENE} the set of scene identifiers.  Both carriers are bounded by
ProB's \texttt{DEFAULT\_SETSIZE}, so no explicit cardinality constant is
needed.  Two of each suffice: one frame the user closes and one bystander,
one scene that is pushed again and one that arrives new.

% ============================================================
\section{Free Types}
% ============================================================

\begin{zed}
  VIS ::= vOpen | vMinimized | vClosed
\end{zed}

A frame's visibility, and the whole of it.  \texttt{vOpen}: painted as an
inner window.  \texttt{vMinimized}: not painted, carrying a dock pill.
\texttt{vClosed}: not painted and carrying no pill, reachable only by a
named user gesture.  The third value is what this design adds; today
\texttt{vClosed} is not a state but the absence of the frame.

% ============================================================
\section{State}
% ============================================================

The state is flat: one schema, nine components.  Four are the Display's
replica of client-owned content, three are the Display-owned visibility, and
two are ghosts --- history variables that record what the \emph{user} did, so
that a later state can be judged against it.  The ghosts are a modelling
device with no counterpart in the code, and they are named where they appear.

\begin{schema}{Sys}
  frames : \power FRAME \\
  content : \power SCENE \\
  sceneFrame : SCENE \pfun FRAME \\
  vis : FRAME \pfun VIS \\
  activeTab : FRAME \pfun SCENE \\
  focus : \power FRAME \\
  userClosed : \power FRAME \\
  closedScenes : \power SCENE \\
  autoFocused : \power FRAME \\
  tabStolen : \power FRAME
\where
  \dom vis = frames \\
  \dom activeTab = frames \\
  \dom sceneFrame = content \\
  \ran sceneFrame \subseteq frames \\
  \forall f : frames @ f \in \ran sceneFrame \\
  \forall f : frames @ activeTab~f \in content \\
  \forall f : frames @ sceneFrame~(activeTab~f) = f \\
  focus \subseteq frames \\
  \# focus \leq 1
\end{schema}

\texttt{content} is the set of scene ids the replica holds --- exactly the
domain of \texttt{sceneFrame}, since every scene lives in a frame.  Membership
of \texttt{content} is what the code computes as \texttt{is\_new}: a scene id
outside it is \emph{unseen}, one inside it is \emph{known}.

\texttt{focus} is the one-shot focus request the renderer consumes, modelled
as a set of size at most one so the empty set can stand for ``none'' without
a sentinel.

The predicate carries only \emph{structural} invariants: the coverage of the
three partial functions, the no-husk rule (every frame holds at least one
scene, which is why \texttt{frame\_id} and its content appear and disappear
together), the active tab pointing at a scene of its own frame, and the
cardinality of the focus request.  These hold in every design, correct or
buggy, so they act as coherence constraints rather than as the properties
under test.  The safety properties (Section~\ref{sec:inv}) are deliberately
\emph{absent} here: a predicate placed in \texttt{Sys} is enforced as a guard
on every successor state, which would silently disable any operation that
broke it and mask the very defects this model exists to catch.

% ============================================================
\section{Initialization}
% ============================================================

\begin{schema}{Init}
  Sys'
\where
  frames' = \emptyset \\
  content' = \emptyset \\
  sceneFrame' = \emptyset \\
  vis' = \emptyset \\
  activeTab' = \emptyset \\
  focus' = \emptyset \\
  userClosed' = \emptyset \\
  closedScenes' = \emptyset \\
  autoFocused' = \emptyset \\
  tabStolen' = \emptyset
\end{schema}

A single initial state: a Display holding nothing.

% ============================================================
\section{Content Operations}
% ============================================================

The three operations a client can cause.  Together they are the whole of
\texttt{SceneReplica.handle\_framed\_scene}.  Not one of them writes
\texttt{vis}, \texttt{focus}, or the \texttt{activeTab} of a frame that
already has one --- that is the R8 rule, and stating it operation by
operation is what makes it checkable.

% -------------------------------------------------------------------
\subsection{Push a New Scene into a New Frame}
% -------------------------------------------------------------------

An unseen scene naming a frame that does not exist yet.  This is the one
place a content event sets a visibility value, and it is sound precisely
because there is no prior value to override: a frame is born
\texttt{vOpen}.  That choice is the \emph{Display's} new-frame policy, not
something the Hub forces --- the Hub sends a scene, and the Display decides
what a frame containing it looks like when it first exists.  Being born open
is not a raise: no focus is requested, no other frame is disturbed, and the
Display window's own shown/hidden state (R1) is untouched.

\begin{schema}{PushNewFrame}
  \Delta Sys \\
  s? : SCENE \\
  f? : FRAME
\where
  s? \notin content \\
  f? \notin frames \\
  frames' = frames \cup \{ f? \} \\
  content' = content \cup \{ s? \} \\
  sceneFrame' = sceneFrame \cup \{ s? \mapsto f? \} \\
  vis' = vis \cup \{ f? \mapsto vOpen \} \\
  activeTab' = activeTab \cup \{ f? \mapsto s? \} \\
  focus' = focus \\
  userClosed' = userClosed \\
  closedScenes' = closedScenes \\
  autoFocused' = autoFocused \\
  tabStolen' = tabStolen
\end{schema}

% -------------------------------------------------------------------
\subsection{Push a New Scene into an Existing Frame}
% -------------------------------------------------------------------

An unseen scene joining a frame that is already there --- a second tab.  The
frame's visibility is whatever the user last made it, including
\texttt{vClosed}, and stays that way; the new scene joins the tab strip
without taking the selection.

\begin{schema}{PushNewScene}
  \Delta Sys \\
  s? : SCENE \\
  f? : FRAME
\where
  s? \notin content \\
  f? \in frames \\
  content' = content \cup \{ s? \} \\
  sceneFrame' = sceneFrame \cup \{ s? \mapsto f? \} \\
  frames' = frames \\
  vis' = vis \\
  activeTab' = activeTab \\
  focus' = focus \\
  userClosed' = userClosed \\
  closedScenes' = closedScenes \\
  autoFocused' = autoFocused \\
  tabStolen' = tabStolen
\end{schema}

% -------------------------------------------------------------------
\subsection{Push a Scene Already Held}
% -------------------------------------------------------------------

The repaint: \texttt{update()}, a beads reload, a track change.  It replaces
element content the Display renders, which this model does not carry, and it
changes nothing else at all.  The $\Xi$ says the whole rule.

\begin{schema}{PushRepeat}
  \Xi Sys \\
  s? : SCENE
\where
  s? \in content
\end{schema}

A scene moving from one frame to another --- the branch of
\texttt{upsert\_scene\_in\_frame} that dismisses a scene from its old frame
first --- is deliberately not modelled.  It is a content-to-content
relocation with no visibility consequence, and adding it would grow the state
space without touching the properties under test.

% -------------------------------------------------------------------
\subsection{Dispose a Scene}
% -------------------------------------------------------------------

The one way ``known'' is legitimately lost: an empty push, a Hub manifest
that disowns the scene, or a frame TTL that has passed.  The client is saying
the content is gone, and after it the id is genuinely unseen again.

When other scenes remain in the frame, the frame stays and the active tab
falls back to a survivor, named by \texttt{keep?}:

\begin{schema}{DisposeScene}
  \Delta Sys \\
  s? : SCENE \\
  f? : FRAME \\
  keep? : SCENE
\where
  s? \in content \\
  sceneFrame~s? = f? \\
  keep? \in \dom (sceneFrame \rres \{ f? \}) \\
  keep? \neq s? \\
  content' = content \setminus \{ s? \} \\
  sceneFrame' = \{ s? \} \ndres sceneFrame \\
  activeTab' = activeTab \oplus \{ f? \mapsto keep? \} \\
  frames' = frames \\
  vis' = vis \\
  focus' = focus \\
  userClosed' = userClosed \\
  closedScenes' = closedScenes \setminus \{ s? \} \\
  autoFocused' = autoFocused \\
  tabStolen' = tabStolen
\end{schema}

When the disposed scene was the frame's last, the frame goes with it --- a
frame with no content is a husk, and that holds however the user had left it,
so a \emph{closed} frame whose content is disposed is disposed too.  The
frame id is released from the ghosts along with everything else: the user's
close applied to a frame that no longer exists, and the next arrival under
that id is a genuinely new frame.

\begin{schema}{DisposeFrame}
  \Delta Sys \\
  s? : SCENE \\
  f? : FRAME
\where
  s? \in content \\
  sceneFrame~s? = f? \\
  \dom (sceneFrame \rres \{ f? \}) = \{ s? \} \\
  content' = content \setminus \{ s? \} \\
  sceneFrame' = \{ s? \} \ndres sceneFrame \\
  frames' = frames \setminus \{ f? \} \\
  vis' = \{ f? \} \ndres vis \\
  activeTab' = \{ f? \} \ndres activeTab \\
  focus' = focus \setminus \{ f? \} \\
  userClosed' = userClosed \setminus \{ f? \} \\
  closedScenes' = closedScenes \setminus \{ s? \} \\
  autoFocused' = autoFocused \setminus \{ f? \} \\
  tabStolen' = tabStolen \setminus \{ f? \}
\end{schema}

% ============================================================
\section{Visibility Operations}
% ============================================================

The operations a user can cause.  Not one of them writes \texttt{content} or
\texttt{sceneFrame}: putting a window away is not throwing its contents out.

% -------------------------------------------------------------------
\subsection{Minimize}
% -------------------------------------------------------------------

\begin{schema}{Minimize}
  \Delta Sys \\
  f? : FRAME
\where
  f? \in frames \\
  vis~f? = vOpen \\
  vis' = vis \oplus \{ f? \mapsto vMinimized \} \\
  focus' = focus \setminus \{ f? \} \\
  frames' = frames \\
  content' = content \\
  sceneFrame' = sceneFrame \\
  activeTab' = activeTab \\
  userClosed' = userClosed \\
  closedScenes' = closedScenes \\
  autoFocused' = autoFocused \\
  tabStolen' = tabStolen
\end{schema}

% -------------------------------------------------------------------
\subsection{Close}
% -------------------------------------------------------------------

The user clicks a frame's close button.  It is a visibility transition and
nothing else: the frame stays in \texttt{frames} so a later raise has
something to act on, and its scenes stay in \texttt{content} so a later push
still reads as a repeat.  \texttt{userClosed} is the ghost recording that the
user, not the system, put this frame away; it is cleared only by a raise or
by the frame's disposal.

\begin{schema}{Close}
  \Delta Sys \\
  f? : FRAME
\where
  f? \in frames \\
  vis~f? \in \{ vOpen, vMinimized \} \\
  vis' = vis \oplus \{ f? \mapsto vClosed \} \\
  focus' = focus \setminus \{ f? \} \\
  userClosed' = userClosed \cup \{ f? \} \\
  closedScenes' = closedScenes \cup \dom (sceneFrame \rres \{ f? \}) \\
  frames' = frames \\
  content' = content \\
  sceneFrame' = sceneFrame \\
  activeTab' = activeTab \\
  autoFocused' = autoFocused \\
  tabStolen' = tabStolen
\end{schema}

% -------------------------------------------------------------------
\subsection{Set the Dock State}
% -------------------------------------------------------------------

A client adjusting a frame's dock state through \texttt{set\_frame\_state} ---
the operation the whole client surface carries (MCP, REST, CLI).  It is not a
user gesture, so it is \emph{not} a way back for a closed frame: its guard
admits only a frame that is up or docked, and it may set only those two values.

Both halves of that guard matter, and for the same reason.  Undocking a closed
frame would put back a window the user shut, with nothing behind it but a client
holding a frame id; docking one would give it a pill in the bar the user never
asked for.  They are one invariant broken in two directions, and invariant~4
catches either.

This operation was missing from the first version of this model, which is why
the model did not catch the defect: the code retained the closed frame (so the
call no longer failed with ``not found'') and then restored it unguarded.  A
specification that omits an operation cannot constrain it.

\begin{schema}{SetDockState}
  \Delta Sys \\
  f? : FRAME \\
  v? : VIS
\where
  f? \in frames \\
  vis~f? \in \{ vOpen, vMinimized \} \\
  v? \in \{ vOpen, vMinimized \} \\
  vis' = vis \oplus \{ f? \mapsto v? \} \\
  frames' = frames \\
  content' = content \\
  sceneFrame' = sceneFrame \\
  activeTab' = activeTab \\
  focus' = focus \\
  userClosed' = userClosed \\
  closedScenes' = closedScenes \\
  autoFocused' = autoFocused \\
  tabStolen' = tabStolen
\end{schema}

It leaves \texttt{focus} alone rather than clearing it, matching the code: only
a close releases a pending focus request, and a bulk dock does not.  It never
writes \texttt{userClosed}, which is what keeps invariant~4 true --- a frame in
that ghost has \texttt{vis} of \texttt{vClosed}, so the guard excludes it and no
run of this operation can reach a state where the user's close still stands over
an open frame.

% -------------------------------------------------------------------
\subsection{Raise}
% -------------------------------------------------------------------

The whole gesture behind a user asking for a frame by name: an applet's menu
entry click (DES-063's raise-first pattern), a dock pill, Expand All, the
Windows menu's list of closed frames.  Restoring and focusing are one
operation, not two --- a frame that took focus while still put away would
have answered the request without becoming visible.  It is enabled from
\emph{any} visibility, which is what makes a closed frame reachable again,
and it is the only operation that both opens a frame and takes focus.

\begin{schema}{Raise}
  \Delta Sys \\
  f? : FRAME
\where
  f? \in frames \\
  vis' = vis \oplus \{ f? \mapsto vOpen \} \\
  focus' = \{ f? \} \\
  userClosed' = userClosed \setminus \{ f? \} \\
  closedScenes' = closedScenes \setminus \dom (sceneFrame \rres \{ f? \}) \\
  frames' = frames \\
  content' = content \\
  sceneFrame' = sceneFrame \\
  activeTab' = activeTab \\
  autoFocused' = autoFocused \\
  tabStolen' = tabStolen
\end{schema}

% -------------------------------------------------------------------
\subsection{Select a Tab}
% -------------------------------------------------------------------

The user clicks a tab in a multi-scene frame.  It is here so that a tab
selection has a user-owned source, against which a content-driven one can be
told apart.

\begin{schema}{SelectTab}
  \Delta Sys \\
  f? : FRAME \\
  s? : SCENE
\where
  f? \in frames \\
  s? \in \dom (sceneFrame \rres \{ f? \}) \\
  activeTab' = activeTab \oplus \{ f? \mapsto s? \} \\
  frames' = frames \\
  content' = content \\
  sceneFrame' = sceneFrame \\
  vis' = vis \\
  focus' = focus \\
  userClosed' = userClosed \\
  closedScenes' = closedScenes \\
  autoFocused' = autoFocused \\
  tabStolen' = tabStolen
\end{schema}

% -------------------------------------------------------------------
\subsection{Consume the Focus Request}
% -------------------------------------------------------------------

The renderer asks once, on the render after the request, and the book forgets
it.  Modelled so that focus does not accumulate and the system stays live.

\begin{schema}{ConsumeFocus}
  \Delta Sys \\
  f? : FRAME
\where
  f? \in focus \\
  focus' = \emptyset \\
  frames' = frames \\
  content' = content \\
  sceneFrame' = sceneFrame \\
  vis' = vis \\
  activeTab' = activeTab \\
  userClosed' = userClosed \\
  closedScenes' = closedScenes \\
  autoFocused' = autoFocused \\
  tabStolen' = tabStolen
\end{schema}

% ============================================================
\section{Invariants}
\label{sec:inv}
% ============================================================

Six properties must hold across every reachable state.  The first five are
state predicates; the sixth is a property of the transition relation.  None is
placed in \texttt{Sys}, for the reason given with that schema.

\paragraph{Invariant 1 --- content never takes focus.}
No focus request is ever raised by a content event:
\[
  autoFocused = \emptyset.
\]
The ghost is written by no operation in this specification, so the property is
established by construction here and is a \emph{reachability} question only in
the fidelity control, where the \texttt{is\_new} side effect writes it.  This
is criterion R8 in its narrowest form: \texttt{show()} and \texttt{update()}
never raise, whether or not the scene is new.

\paragraph{Invariant 2 --- content never steals the tab.}
No active-tab change is ever caused by a content arrival:
\[
  tabStolen = \emptyset.
\]
The active tab is a user-owned selection, in the same category as visibility.
A new scene may join a frame's tab strip; it may not pull the selection off
whatever the user was reading.  \texttt{DisposeScene} does repoint the tab,
but only away from a scene that has ceased to exist and only to a survivor,
which is a repair, not a steal.

\paragraph{Invariant 3 --- a closed frame is still there to raise.}
Every frame the user has closed is still a frame:
\[
  userClosed \subseteq frames.
\]
This is bug A stated as a property.  A closed frame that has left
\texttt{frames} cannot satisfy \texttt{Raise}'s precondition, so the user's
gesture --- an applet menu entry, a Windows-menu entry --- has nothing to act
on, and the frame is reachable again only by whatever accident recreates it.

\paragraph{Invariant 4 --- only a raise reopens a closed frame.}
No frame is on screen while the user's close still stands:
\[
  \{ f : frames | f \in userClosed \land vis~f = vOpen \} = \emptyset.
\]
This is bug B stated as a property.  \texttt{Raise} is the only operation that
sets \texttt{vOpen} on an existing frame, and it clears \texttt{userClosed} in
the same step, so the two can never be true of one frame at once.  Any state
where they are is a frame that came back without the user asking.

\paragraph{Invariant 5 --- the content that was put away stays put away.}
No scene the user closed is back on screen while that close still stands:
\[
  \forall s : closedScenes \cap content @ vis~(sceneFrame~s) \neq vOpen.
\]
This is bug B in the shape the user actually meets it: not merely that some
frame id came back, but that the very scene they shut --- voxd's music board ---
is in front of them again on a track change.  Invariant~4 is about frame
identity; this one is about the content inside it, and it is the sharper of the
two because a frame id could in principle be reused for something else.
\texttt{closedScenes} is the ghost holding the scene ids that were in a frame at
the moment it was closed, cleared when that frame is raised or when the scene is
disposed.

\paragraph{Invariant 6 --- deadlock freedom.}
No reachable state leaves every operation disabled.  \texttt{DisposeFrame}
returns a frame id and a scene id to their unseen pools, so a workspace can
always be emptied and repopulated and the system cycles rather than
terminating; and \texttt{PushNewFrame} is enabled in the empty state.  This is
a property of the reachable state space, discharged by the model check in
Section~\ref{sec:verify} rather than by a goal.

% ============================================================
\section{Verification}
\label{sec:verify}
% ============================================================

Type-checking is a \texttt{fuzz} pass:
\begin{verbatim}
  fuzz -t docs/frame-visibility-lifecycle.tex
\end{verbatim}

The four invariants are state properties, checked by asking ProB whether the
\emph{negation} of each is reachable.  A goal that is not found means the
invariant holds on every reachable state:
\begin{verbatim}
  # Invariant 1 (must report: goal NOT found)
  probcli docs/frame-visibility-lifecycle.tex -model_check -nodead \
    -goal "autoFocused /= {}" -p DEFAULT_SETSIZE 2

  # Invariant 2 (must report: goal NOT found)
  probcli docs/frame-visibility-lifecycle.tex -model_check -nodead \
    -goal "tabStolen /= {}" -p DEFAULT_SETSIZE 2

  # Invariant 3 (must report: goal NOT found)
  probcli docs/frame-visibility-lifecycle.tex -model_check -nodead \
    -goal "userClosed /\ frames /= userClosed" -p DEFAULT_SETSIZE 2

  # Invariant 4 (must report: goal NOT found)
  probcli docs/frame-visibility-lifecycle.tex -model_check -nodead \
    -goal "card({f|f:FRAME & f:frames & f:userClosed & vis(f)=vOpen})>0" \
    -p DEFAULT_SETSIZE 2

  # Invariant 5 (must report: goal NOT found)
  probcli docs/frame-visibility-lifecycle.tex -model_check -nodead \
    -goal "#s.(s:SCENE & s:content & s:closedScenes
               & vis(sceneFrame(s))=vOpen)" \
    -p DEFAULT_SETSIZE 2
\end{verbatim}

A goal predicate may not name an element of a given set: \texttt{probcli}
rejects \texttt{FRAME2 : frames} with a \texttt{type\_expression\_error}
even though it prints those names in its own traces, and the model-check then
reports ``no counter example'' \emph{vacuously}.  Every goal above is therefore
written over the state components and bound variables only.  That limitation is
why Invariant~5 needs the \texttt{closedScenes} ghost at all: without it, ``the
same scene came back'' cannot be stated in a form ProB will accept.

Invariant~6 is the deadlock check over the full reachable state space, which
is what \texttt{-model\_check} performs when no \texttt{-nodead} is given:
\begin{verbatim}
  probcli docs/frame-visibility-lifecycle.tex -model_check \
    -p DEFAULT_SETSIZE 2
\end{verbatim}

\paragraph{Carrier sensitivity.}  The results above are at
\texttt{DEFAULT\_SETSIZE 2}.  The model is also deadlock-free and
counter-example-free at \texttt{DEFAULT\_SETSIZE 3} --- 2593 states, 34939
transitions, \texttt{ALL OPERATIONS COVERED} --- which is worth having because
three frames and three scenes is the first size at which a frame can be closed
while two others remain in different visibilities.  Size 4 does not finish in a
useful time and is not required: every operation is already covered at 2, and
the properties are not cardinality-dependent.

\paragraph{\texttt{-cbc\_deadlock} reports a spurious deadlock here, and on
every other Z specification in this repository.  Ignore it.}
The \texttt{make prob} target runs a constraint-based deadlock check as well as
the model check, and that step prints \texttt{*** DEADLOCK state found ***} for
this document.  The state it reports is not a state of this specification:

\begin{verbatim}
  vis = {(FRAME1|->vOpen), (FRAME1|->vMinimized)},  frames = {}
\end{verbatim}

\noindent
\texttt{vis} maps one frame to two values, so it is not a function at all, let
alone the partial function \texttt{Sys} declares; and \texttt{\dom vis} is
\texttt{\{FRAME1\}} while \texttt{frames} is empty, so it violates
\texttt{Sys}'s own \texttt{\dom vis = frames}.  In ProB's Z mode
\texttt{-cbc\_deadlock} does not constrain its solutions to the state schema.
The clearest demonstration is a two-line control --- a counter declared
$n : \nat$ and bounded $n \leq 2$, whose only genuine deadlock is at $n = 2$:
\texttt{-cbc\_deadlock} reports the deadlocking state as \textbf{$n = -2$},
having enumerated $n$ over \texttt{INTEGER}.

Three consequences.  It is not a reachable deadlock: the reachable-state check
says so directly.  It is not an artifact of a small carrier: it persists
unchanged at \texttt{DEFAULT\_SETSIZE} 2, 3, 4 and 6.  And it is not specific
to this document: \texttt{display\_lifecycle.tex},
\texttt{architecture/workspace-model.tex}, \texttt{hub\_replicator.tex} and
\texttt{listen\_lifecycle.tex} all report it too --- including
\texttt{display\_lifecycle.tex}, whose own Verification section discharges
deadlock freedom with \texttt{-model\_check} and is correct to do so.

The reachable-state check is the authoritative one, and its deadlock detection
is live rather than vacuous: run against the $n \leq 2$ control it reports
\texttt{Found "deadlock" error in state id 2}, the real one.

\paragraph{Results.}  Run against ProB 1.15.1 (SICStus 4.8.0).  The full
reachable state space is \textbf{119 states, 1309 transitions}, with
\texttt{ALL OPERATIONS COVERED} --- every one of the eleven operation schemas
fires, so no invariant is discharged vacuously by an operation that never
ran --- and no deadlock.  All five goals report \emph{no counter example
found, all states visited}:

\begin{center}
\begin{tabular}{llr}
  \textbf{Invariant} & \textbf{Goal} & \textbf{Result} \\
  \hline
  1 & \texttt{autoFocused /= \{\}} & not found \\
  2 & \texttt{tabStolen /= \{\}} & not found \\
  3 & \texttt{userClosed /\char`\\ frames /= userClosed} & not found \\
  4 & \texttt{card(\{f|...\ vis(f)=vOpen\})>0} & not found \\
  5 & \texttt{\#s.(s:closedScenes \& vis(sceneFrame(s))=vOpen)} & not found \\
  6 & deadlock (no goal --- the model check itself) & none \\
\end{tabular}
\end{center}

The \texttt{fuzz} leg passes on both this document and the control.  The
fidelity control reaches every one of the five (Section~\ref{sec:fidelity}),
so none of these is a vacuous pass.

\paragraph{\texttt{SetDockState} and the defect that only the new design can
have.}  This operation was added after the first implementation round, when
review found that \texttt{set\_frame\_state} --- carried by the whole client
surface --- restored a closed frame with no gesture behind it.  The model had
not caught it for the plain reason that the model did not have the operation.

The evidence for the guard is a deterministic replay rather than a search-order
witness, for the reason given in Section~\ref{sec:fidelity}:

\begin{verbatim}
  machine('frame-visibility-lifecycle').
  '$initialise_machine'.
  'PushNewFrame'(fd(1,'FRAME'),fd(1,'SCENE')).
  'Close'(fd(1,'FRAME')).
  'SetDockState'(fd(1,'FRAME'),vOpen).
\end{verbatim}

\noindent
Replayed with \texttt{probcli ... -t -p DEFAULT\_SETSIZE 2} it \textbf{fails at
step 3}, naming \texttt{SetDockState} as the operation that is not enabled:
the guard $vis~f? \in \{vOpen, vMinimized\}$ excludes a closed frame.  Remove
that guard and the same sequence replays into a state where $f_1 \in userClosed$
and $vis~f_1 = vOpen$ --- the negation of invariant~4.

The usual fidelity control cannot host this one, and the reason is worth
stating: its \texttt{Close} is the destructive original, which \emph{removes}
the frame, so \texttt{SetDockState}'s $f? \in frames$ fails and the sequence is
not expressible there at all.  That is the point rather than a gap.  This defect
is reachable only in a design that \emph{retains} the closed frame, so it is a
defect the old code could not have and the new code created --- an accident of
the fix, closed here by the guard.  What the control does show is the operation
running unguarded: \texttt{SetDockState}$(f_1, vClosed)$ then
\texttt{SetDockState}$(f_1, vOpen)$ replays successfully there, a client both
shutting a window and reopening it with no user gesture anywhere.

% ============================================================
\section{Fidelity: reproducing bugs A and B}
\label{sec:fidelity}
% ============================================================

A model that cannot reproduce the known defects proves nothing.  The fidelity
control is \texttt{docs/frame-visibility-lifecycle\_buggy.tex}: this
specification with today's behaviour restored in exactly three schemas ---
\texttt{Close} disposes the frame instead of marking it closed, and the two
new-scene pushes carry the \texttt{is\_new} side effect (clear the minimized
flag, request focus, take the tab).  \texttt{Sys}, \texttt{Init},
\texttt{PushRepeat}, \texttt{DisposeScene}, \texttt{DisposeFrame},
\texttt{Minimize}, \texttt{Raise}, \texttt{SelectTab} and
\texttt{ConsumeFocus} are identical.  All five goals above must be
\emph{found} there, and all five are: its state space is 1030 states and
10309 transitions, also with \texttt{ALL OPERATIONS COVERED}.  The traces
below are ProB's, not hand-derived.

\paragraph{A note on witness traces.}  A goal's counter-example trace is
\emph{not} stable across runs: ProB's search order varies, and repeated
invocations of the same goal return witnesses of different lengths.  What is
stable is the verdict (found / not found) and the state-space size.  So the
fidelity evidence below is given as two \emph{trace replays} --- an explicit
sequence of operations checked with \texttt{probcli -t} against a
\texttt{.trace} file --- because a replay is deterministic and names exactly
which operation is enabled in which design.

\paragraph{Bug B --- the closed frame reopens on a background push.}
The sequence is voxd's: install the music scene, the user shuts the window, the
track changes and the same scene is pushed into the same frame.

\begin{verbatim}
  machine('frame-visibility-lifecycle').
  '$initialise_machine'.
  'PushNewFrame'(fd(1,'FRAME'),fd(1,'SCENE')).
  'Close'(fd(1,'FRAME')).
  'PushNewFrame'(fd(1,'FRAME'),fd(1,'SCENE')).
\end{verbatim}

Replayed with \texttt{probcli docs/frame-visibility-lifecycle.tex -t -p
DEFAULT\_SETSIZE 2}, this \textbf{fails at step 3}: \texttt{PushNewFrame} is
not enabled, because \texttt{Close} left $s_1$ in \texttt{content} and the
operation's guard is $s_1 \notin content$.  The push is a \texttt{PushRepeat}
instead --- a $\Xi$ operation --- and the frame stays \texttt{vClosed}.
Against the control (same file name changed to the buggy machine) the identical
sequence \textbf{replays successfully}, and the frame is back on screen and
focused.  That one operation being enabled or not is the whole of bug B, and it
is created by the close, not by the push.

\paragraph{Bug A --- the frame the user closed cannot be raised.}
The same shape, in the opposite direction: install a board, the user shuts it,
the user clicks the applet's menu entry.

\begin{verbatim}
  machine('frame-visibility-lifecycle').
  '$initialise_machine'.
  'PushNewFrame'(fd(1,'FRAME'),fd(1,'SCENE')).
  'Close'(fd(1,'FRAME')).
  'Raise'(fd(1,'FRAME')).
\end{verbatim}

This \textbf{replays successfully} here --- the frame is still in
\texttt{frames}, so \texttt{Raise} is enabled and restores it --- and
\textbf{fails at step 3} against the control, where \texttt{Raise}'s
precondition $f_1 \in frames$ cannot hold because \texttt{Close} removed the
frame.  In the running system that failure is \texttt{raise\_frame} answering
\texttt{raised: false} to an applet's menu click.

\paragraph{The focus and tab steals.}  Both are reachable in the control within
two steps and unreachable here: \texttt{PushNewFrame} alone leaves
\texttt{autoFocused} non-empty there, and a following
\texttt{PushNewScene} into the same frame leaves \texttt{tabStolen} non-empty.
Neither ghost is ever written by any operation in this specification.

\paragraph{What the difference establishes.}  The reachability searches return
\emph{goal not found} for this specification --- all five, over 95 states with
every operation covered --- and \emph{goal found} for the control, all five, on
the same carriers and the same initial state.  The two replays sharpen that into
a statement about single operations: the step that reproduces each bug in the
control is \emph{not enabled} here, and the step that repairs it here is not
enabled there.  The model is therefore not too abstract to see the defects: it
distinguishes exactly the designs that exhibit them from the design that does
not.

% ============================================================
\section{What this model does not carry}
% ============================================================

Three things are deliberately out of scope, and are named so that a reader
does not mistake silence for a claim.

\begin{itemize}
  \item \textbf{The Display window's own shown/hidden state} (R1 of the same
    epic).  This model is about frames \emph{within} the workspace.  The rule
    it establishes --- a content event never changes visibility --- is the
    same rule R1 applies one level up, and nothing here contradicts it: a
    frame born \texttt{vOpen} inside a hidden window is still invisible.
  \item \textbf{Persistence across a Display restart.}  \texttt{vClosed}, like
    \texttt{vMinimized}, focus, and cascade position, is Display-local session
    state; a restarted Display rebuilds its replica from the Hub's manifest
    and every frame is created fresh.  Making the close survive a restart
    would require replicating a Display-owned value back to the Hub, which is
    the coupling this design exists to sever.
  \item \textbf{Scene content and element identity.}  A push carries an
    element tree, and replacing one retires the element ids it dropped.  That
    is modelled in \texttt{docs/patch\_application.tex} and is orthogonal to
    visibility.
  \item \textbf{Resolving a caller's name to a store key.}  \texttt{Raise}
    takes \texttt{f? : FRAME} as already a well-formed member of the
    carrier --- the schema has no notion of a caller naming a frame by a
    plain local id distinct from the connection-scoped key the Display
    actually stores it under.  That resolution (a Hub-side registry lookup,
    not a Display-side visibility transition) is a missing method call, not
    a new interleaving: no operation here reads or writes the mapping from a
    caller's name to \texttt{FRAME}, so a defect in that mapping cannot
    surface as a violated invariant of this model, and fixing it changes no
    schema above.  The empirical regression coverage for that resolution
    lives beside the fix, not here.
\end{itemize}

\end{document}
